\section{Fórmulas de Reiner y Rubinstein}
\label{sec:impl_rr}

% OBJETIVO: Documentar tu implementación de las fórmulas de R&R
% Este es tu MÉTODO 1

\subsection{Descripción de la Implementación}
\label{subsec:desc_impl_rr}

% QUÉ ESCRIBIR AQUÍ:
% - Cómo implementaste las fórmulas
% - Lenguaje y librerías
% - Estructura del código

\subsubsection{Cálculo de Variables Auxiliares}

% Explicar:
% - Implementación de r, d, σ, t, φ, η
% - Cálculo de μ, λ, x, x_1, y, y_1, z, b

\subsubsection{Cálculo de Términos de Valoración}

% Explicar:
% - Implementación de términos [1] a [6]
% - Uso de funciones de densidad f(u), g(u), h(τ)
% - Cálculo de N(·) (CDF normal)

\subsubsection{Construcción de Fórmulas para Cada Tipo}

% Explicar:
% - Cómo implementas las 8 fórmulas (DI, DO, UI, UO para Call y Put)
% - Manejo de casos especiales (K > H, K < H)
% - Uso de relaciones de paridad

\subsection{Detalles de Implementación}
\label{subsec:detalles_impl_rr}

% QUÉ ESCRIBIR AQUÍ:
% - Lenguaje: Python/C++/MATLAB/etc
% - Librerías: NumPy, SciPy, etc
% - Snippets de código

\subsubsection{Código Relevante}

% INCLUIR CÓDIGO:
% Ejemplo de función para calcular Down-and-Out Call

\subsection{Cálculo de Griegas}
\label{subsec:griegas_rr}

% SI CALCULAS GRIEGAS:
% - Delta, Gamma, Vega, Theta
% - Cómo las calculas (derivadas analíticas o diferencias finitas)

\subsection{Validación}
\label{subsec:validacion_rr}

% QUÉ ESCRIBIR AQUÍ:
% - Verificar propiedades matemáticas
% - Verificar relaciones de paridad
% - Verificar límites asintóticos

\subsubsection{Verificación de Relaciones de Paridad}

% Verificar:
% - C_DI + C_DO = C_vanilla
% - C_UI + C_UO = C_vanilla
% - Similar para puts

\subsubsection{Verificación de Límites Asintóticos}

% Verificar:
% - Cuando H → 0 (barrera muy baja): C_DO → C_vanilla
% - Cuando H → ∞ (barrera muy alta): C_UO → C_vanilla
% - Cuando σ → 0: convergencia a payoffs determinísticos

\subsubsection{Ejemplos Numéricos}

% INCLUIR 3-4 EJEMPLOS:
% - Down-and-Out Call
% - Down-and-In Call
% - Up-and-Out Put
% - Verificar que paridad se cumple
